% from ega2-errata-addenda-fr.tex
\oldpage[II]{217}
\markboth{ERRATA ET ADDENDA, LISTE 1}{ERRATA ET ADDENDA, LISTE 1}
\phantomsection
\addcontentsline{toc}{section}{Errata et Addenda (Liste 1)}
\section*{Errata et Addenda}
\thispagestyle{plain}
\begin{center}
\large (Liste 1)
\end{center}

\subsection*{Chapitre 0}

\paragraph{(0, 2.1.3).}
Ligne 15 du bas de la p.~21, remplacer «~inter-section~» par
«~intersection~».

\paragraph{(0, 2.1.6).}
Ligne 13 de la p.~22, remplacer
\[
  Z_i\cap C\left(\bigcup_{j\ne i}Z_j\right)
  \quad\hbox{par}\quad
  \bigcup_{j\ne i}Z_j.
\]

\paragraph{(0, 3.1.6).}
Ligne 17 du bas de la p.~25, remplacer
$\Gamma(U,\sh{F})$ par $\Gamma(V,\sh{F})$.

\paragraph{(0, 4.1.1).}
Ligne 12 de la p.~36, remplacer $\sh{G}\to\sh{F}$ par
$\sh{B}\to\sh{A}$.

\paragraph{(0, 4.2.1).}
Ligne 11 du bas de la p.~39, remplacer «~faisceaux~» par «~faisceau~».

\paragraph{(0, 5.1.2).}
Après la ligne 14 de la p.~45, ajouter :

Si $\sh{G}$ est un $\sh{O}_Y$-Module engendré par ses sections au-dessus
de $Y$, alors $f^*(\sh{G})$ est engendré par ses sections au-dessus de
$X$, car $f^*$ est un foncteur exact à droite.

\paragraph{(0, 5.2.4).}
Ligne 18 du bas de la p.~46, remplacer $f_U$ par $f_U^*$.

\paragraph{(0, 5.3.4).}
Avant la ligne 16 du bas de la p.~47, ajouter :

Si
\[
  \sh{A}\to\sh{B}\to\sh{C}\to\sh{D}\to\sh{E}
\]
est une suite exacte de $\sh{O}_X$-Modules, et si
$\sh{A},\sh{B},\sh{D},\sh{E}$ sont cohérents, alors $\sh{C}$ est
cohérent.

\paragraph{(0, 5.4.1).}
Après la ligne 2 de la p.~49, ajouter :

Supposons $\sh{O}_X$ cohérent, et soit $\sh{F}$ un $\sh{O}_X$-Module
cohérent. Alors, si en un point $x\in X$, $\sh{F}_x$ est un
$\sh{O}_x$-module libre de rang $n$, il existe un voisinage $U$ de $x$
tel que $\sh{F}|U$ soit localement libre de rang $n$; en effet
$\sh{F}_x$ est alors isomorphe à $\sh{O}_x^n$, et la proposition résulte
de (0, 5.2.7).

\paragraph{(0, 6.6.2).}
Remplacer les lignes 10 et 11 de la p.~59 par :

En effet, cela résulte de (0, 6.4.1, d)).

\paragraph{(0, 6.7.1).}
Avant la ligne 10 du bas de la p.~59, ajouter :

Si $f$ est un morphisme plat, on dit encore que $X$ est \emph{plat sur
$Y$}, ou \emph{$Y$-plat}.

\paragraph{(0, 7.2.9).}
Ligne 11 du bas de la p.~65, remplacer $M_n$ par $M_{n-1}$.

\subsection*{Chapitre I}

\paragraph{(I, 1.2.1).}
Ligne 18 de la p.~83, remplacer (deux fois) $K$ par $k$.

\paragraph{(I, 1.2.7).}
Ligne 16 du bas de la p.~84, remplacer $A$ par $A'$ (deux fois) et
$\mathfrak N$ par $\mathfrak N'$; ligne 17 du bas, remplacer $X$ par
$X'$.

\paragraph{(I, 1.7.3) et (I, 2.2.4).}
Les énoncés donnés dans ces numéros peuvent (selon une remarque due à
J.~Tate) être généralisés comme suit :

\subsection*{1.8. Morphismes d'espaces annelés en anneaux locaux dans les schémas affines}
\label{ega2:addendum:subsection:I.1.8-fr}

\begin{proposition}[1.8.1]
\label{ega2:addendum:I.1.8.1-fr}
Soient $(S,\sh{O}_S)$ un schéma affine, $(X,\sh{O}_X)$ un espace annelé
en anneaux locaux. Il y a alors une bijection canonique de l'ensemble
des homomorphismes d'anneaux
